\chapter{Myositis}

\section*{Definition}
Myositis is inflammation of skeletal muscle, presenting with symmetric proximal muscle weakness, elevated serum creatine kinase (CK), and inflammatory infiltrates on muscle biopsy. Idiopathic inflammatory myopathies include polymyositis, dermatomyositis, and inclusion body myositis.

\section*{Pathophysiology}
In polymyositis, cytotoxic CD8+ T cells invade muscle fibers expressing MHC class I molecules, causing direct myocyte necrosis. Dermatomyositis is a complement-mediated microangiopathy targeting endomysial capillaries, resulting in perifascicular atrophy. Inclusion body myositis involves both inflammatory and degenerative mechanisms with rimmed vacuoles and protein aggregates (TDP-43, beta-amyloid). Elevated intramuscular pressure from edema may further compromise perfusion.

\section*{Etiology}
\begin{itemize}
  \item \textbf{Idiopathic inflammatory myopathy:} Polymyositis, dermatomyositis, inclusion body myositis
  \item \textbf{Infectious:} Viral (influenza, HIV, Coxsackie, hepatitis B/C), bacterial (pyomyositis caused by \textit{S. aureus}), parasitic (trichinosis, toxoplasmosis)
  \item \textbf{Drug-induced:} Statins, hydroxychloroquine, penicillamine, procainamide, TNF inhibitors, checkpoint inhibitors
  \item \textbf{Autoimmune:} Overlap syndromes with scleroderma, SLE, mixed connective tissue disease, Sj\"{o}gren syndrome
  \item \textbf{Paraneoplastic:} Dermatomyositis (especially ovarian, lung, pancreatic, gastric, colorectal cancers)
  \item \textbf{Traumatic:} Crush injury, compartment syndrome, rhabdomyolysis
\end{itemize}

\section*{Indications for Evaluation}
Seek care if:
\begin{itemize}
  \item Progressive symmetric proximal weakness (difficulty rising from chair, climbing stairs, lifting arms)
  \item Muscle pain or tenderness with weakness
  \item Heliotrope rash (purple periorbital) or Gottron papules (extensor finger erythema)
  \item Dark urine (myoglobinuria) suggesting rhabdomyolysis
  \item Dysphagia, shortness of breath, or aspiration symptoms
\end{itemize}

\section*{Related Symptoms}
\begin{itemize}
  \item Proximal muscle weakness (shoulder and hip girdles)
  \item Myalgia and muscle tenderness
  \item Heliotrope rash with periorbital edema (dermatomyositis)
  \item Gottron papules over metacarpophalangeal and interphalangeal joints
  \item Dysphagia, dysphonia, and respiratory muscle weakness
  \item Mechanic's hands (cracked, hyperkeratotic fingertips)
\end{itemize}
\textbf{Note:} Dermatomyositis in adults has a strong association with underlying malignancy; age-appropriate cancer screening is essential at diagnosis.

\section*{Self-Care / Home Management}
\begin{itemize}
  \item Energy conservation and pacing of activities
  \item Modified exercise program: low-impact stretching and gentle resistance
  \item Sun protection (broad-spectrum SPF 50) for photosensitive dermatomyositis rash
  \item Dietary modifications for dysphagia (soft foods, thickened liquids)
  \item Fall prevention strategies within the home
\end{itemize}

\section*{Diagnostic Approach}
\begin{itemize}
  \item Serum CK, aldolase, AST, ALT, LDH, and myoglobin levels
  \item Autoantibody panel: ANA, anti-Jo-1 (anti-synthetase), anti-SRP, anti-MDA5, anti-Mi-2, anti-TIF-1$\gamma$
  \item Electromyography showing myopathic motor units with spontaneous fibrillation potentials
  \item Muscle MRI with STIR sequence for edema and inflammation
  \item Muscle biopsy (vastus lateralis or deltoid) for histopathologic confirmation
  \item Malignancy screening: CT chest/abdomen/pelvis, mammography, pelvic ultrasound, colonoscopy as indicated
\end{itemize}

\section*{Treatment / Management Overview}
\begin{itemize}
  \item High-dose corticosteroids (prednisone 0.5--1 mg/kg/day) as first-line therapy
  \item Steroid-sparing immunosuppressants: methotrexate, azathioprine, mycophenolate mofetil
  \item Intravenous immunoglobulin (IVIg) for refractory dermatomyositis and inclusion body myositis
  \item Rituximab or other biologics for treatment-resistant disease
  \item Physical therapy during disease remission for functional recovery
  \item Speech and swallowing therapy for oropharyngeal involvement
\end{itemize}

\clearpage
